%! Author = murfe
%! Date = 08.06.2026

\chapter{Grundlagen}\label{ch:grundlagen}
\textcolor{red}{Orientierung ca. 8-12 Seiten}
%TODO: Einleitung
%TODO: auf jeden Fall irgendwo at most once / at least once / effectively exactly once aufgreifen

\section{Grundlagen der Hintergrundverarbeitung}\label{sec:hintergrundverarbeitung}
Hintergrundverarbeitung adressiert eine zentrale Herausforderung moderner Backend-Systeme: In verteilten Umgebungen sind Teilausfälle,
Timeouts und kurzzeitige Unterbrechungen üblich, sodass nicht ohne entsprechende Mechanismen davon ausgegangen werden kann, dass eine angestoßene
Verarbeitung vollständig und genau einmal ausgeführt wird.
Insbesondere kann bei Kommunikationsabbrüchen oder Wiederholungsversuchen unklar bleiben, ob eine Operation bereits verarbeitet wurde oder ob sie
erneut angestoßen werden muss \parencite[Kap.~8]{kleppmann_designing_2017}.
Vor diesem Hintergrund werden in den folgenden Abschnitten grundlegende Konzepte vorgestellt, die eine Entkopplung der Verarbeitung ermöglichen und die Basis
für robuste Hintergrundverarbeitung bilden.

\subsection{Synchrone vs. Asynchrone Verarbeitung}\label{subsec:synchroneVsAsynchron}
In Backend-Systemen lässt sich Kommunikation und Verarbeitung bezüglich der zeitlichen Kopplung zwischen synchronem und asynchronem Verhalten unterscheiden:
Bei synchroner Verarbeitung nach dem Request-Response-Stil wartet der Aufrufer auf eine Antwort auf die Anfrage und blockiert gegebenenfalls bevor er fortfährt.
Bei asynchroner Verarbeitung wird die Bearbeitung des Auftrags vom Zeitpunkt des Erhalts entkoppelt.
Der Aufrufer initiiert die Arbeit, blockiert jedoch nicht und kann auch während des Wartens auf eine Antwort weiter arbeiten \parencite[Kap.~3.1.1]{richardson_microservices_2018}.

Um Nachrichtenverluste zu reduzieren und Lastspitzen auszugleichen, verwenden asynchrone Systeme in der Praxis häufig einen Puffer.
%TODO: ggfs. Abbildung zu funktionsweise eines Message Brokers
Durch das sogenannte \glqq queueing\grqq \parencite[Kap.~11, Abschnitt \glqq Message brokers\grqq]{kleppmann_designing_2017}, bei dem erhaltene Anfragen zur späteren Ausführung
in den Puffer geschrieben werden, wird unter anderem die Asynchronität erreicht.
Ein Produzent sendet eine Nachricht an einen Message Broker und wartet in der Regel nur auf das Acknowledgment, dass die Nachricht zwischengespeichert wurde,
nicht jedoch auf die Verarbeitung oder das Ergebnis \parencite[Kap.~11, Abschnitt \glqq Messaging Systems\grqq]{kleppmann_designing_2017}.
Die so erreichte Entkopplung hat gleichzeitig den Vorteil, dass die Verarbeitung der Anfrage unabhängig von der Verbindung zwischen Produzent und Konsument wird.
Im Gegensatz dazu kann bei synchronen Systemen bei Verbindungsabbrüchen oder Timeouts unklar bleiben, ob eine Anfrage den Server erreicht hat oder bereits verarbeitet wurde.
Genau dieses Problem adressieren Message Broker:
Falls ein Konsument temporär nicht verfügbar ist oder während der Verarbeitung ausfällt, bleibt das Acknowledgment aus, sodass unklar ist, ob die Nachricht erfolgreich verarbeitet wurde.
Infolgedessen kann die Zustellung zu einem späteren Zeitpunkt oder an einen anderen Konsumenten erneut versucht werden.
Daher ist die Nutzung von Acknowledgements wichtig, damit der Broker erkennen kann, ob eine Nachricht als erfolgreich verarbeitet gilt oder erneut zugestellt werden muss
\parencite[Kap.~11, Abschnitt \glqq Acknowledgments and redelivery\grqq]{kleppmann_designing_2017}.
Auf die Zustellungssemantik wird in Kapitel ... näher eingegangen.%TODO: Referenz auf Kapitel wo darauf näher eingangen wird
Die synchrone Request-Response-Kommunikation bietet demgegenüber eine unmittelbare Rückmeldung über Erfolg oder Fehler einer Operation, setzt jedoch voraus, dass Client und Server 
während der gesamten Interaktion verfügbar sind \parencite[Kap.~3.4]{richardson_microservices_2018}.

Die asynchrone Verarbeitung bildet damit die Grundlage für die Hintergrundverarbeitung, indem die Entkopplung genutzt wird, um Aufgaben unabhängig vom ursprünglichen Aufruf 
auszuführen und gleichzeitig die Robustheit und Skalierbarkeit des Systems zu erhöhen \parencite[Kap.~8, Abschnitt \glqq Why Use A Message Broker?\grqq]{newman_building_2026}.

\subsection{Background Job-Processing}\label{subsec:backgroundJobs}
Background Job-Processing bezeichnet die Ausführung von Aufgaben (Jobs), die als eigenständige Arbeitseinheiten modelliert und außerhalb des
unmittelbaren Request-Response-Ablaufs verarbeitet werden.[Kleppmann 2017, Kap.8–9]
Solche Jobs treten in Backend-Systemen typischerweise dann auf, wenn Verarbeitung zeitintensiv ist (z.B. Report-Generierung oder Datenimporte),
wenn Abläufe entkoppelt werden sollen oder wenn eine robuste Abarbeitung trotz Teilausfällen erforderlich ist.[Kleppmann 2017, Kap.8,10]
Charakteristisch ist, dass die Bearbeitung nicht an die Lebensdauer eines einzelnen Requests gebunden ist, sondern durch dedizierte Worker-Komponenten erfolgt,
die Jobs abarbeiten.[Kleppmann 2017, Kap.8–9]


\subsubsection*{Architekturmodell und zentrale Komponenten}
Ein verbreitetes Architekturmodell für Background Job-Processing ist die Trennung in (1) eine Komponente, die Jobs erzeugt bzw. einplant (z.B. eine API oder ein Scheduler),
(2) einen persistenten Job-Store und (3) Worker, die Jobs aus dem Store beziehen und ausführen.[Kleppmann 2017, Kap.8–9]
Die Persistenz des Job-Store ist hierbei wesentlich, da in verteilten Systemen jederzeit Teilausfälle auftreten können: Prozesse können unerwartet terminieren,
Maschinen können ausfallen oder temporäre Infrastrukturprobleme können den Zugriff auf Ressourcen verhindern.[Kleppmann 2017, Kap.8]
Ein rein flüchtiges In-Memory-Queueing wäre in solchen Situationen nicht ausreichend, weil Jobs verloren gehen oder ihr Bearbeitungszustand nicht mehr rekonstruierbar wäre.[Kleppmann 2017, Kap.8]
Aus der Notwendigkeit von Persistenz ergibt sich zudem, dass Job-Processing-Systeme üblicherweise ein explizites Zustandsmodell benötigen (z.B. „queued“, „running“, „succeeded“, „failed“),
um den Fortschritt und den aktuellen Bearbeitungsstand zu repräsentieren.[Kleppmann 2017, Kap.9]
Ein solches Statusmodell ist nicht nur für die Laufzeitsteuerung wichtig, sondern bildet auch die Grundlage für Korrektheitsargumente und spätere Nachweise, da sich viele
gewünschte Zusicherungen als Eigenschaften zulässiger Statusübergänge formulieren lassen.[Kleppmann 2017, Kap.9]


\subsubsection*{Parallelität, Koordination und Mehrfachverarbeitung}
Um Durchsatz und Latenz zu verbessern, werden Jobs häufig parallel durch mehrere Worker-Instanzen verarbeitet.[Kleppmann 2017, Kap.9]
Diese Parallelität führt jedoch zu Koordinationsanforderungen: Insbesondere muss verhindert werden, dass mehrere Worker denselben Job gleichzeitig bearbeiten oder
widersprüchliche Statusänderungen vornehmen.[Kleppmann 2017, Kap.9] Daneben muss das System nach einem Prozessabbruch unterscheiden können, ob ein Job tatsächlich
noch „in Bearbeitung“ ist oder ob er lediglich in einem Zwischenzustand verblieben ist, weil der Worker während der Bearbeitung ausgefallen ist.[Kleppmann 2017, Kap.8–9]
Job-Processing-Systeme adressieren diese Anforderungen typischerweise durch Mechanismen, die sicherstellen, dass jeweils genau ein Worker das Recht zur Bearbeitung eines Jobs
erhält (z.B. durch Sperr- oder Claiming-Strategien) und dass dieses Bearbeitungsrecht im Fehlerfall wieder freigegeben werden kann.[Kleppmann 2017, Kap.9]
Die konkrete Ausprägung hängt stark vom gewählten Job-Store ab, etwa ob eine Datenbank als persistenter Store verwendet wird oder ob ein dediziertes Queueing-System
eingesetzt wird.[Kleppmann 2017, Kap.8–9]

\subsubsection*{Fehlerbehandlung und Wiederaufnahme}
Ein weiterer Grundpfeiler von Background Job-Processing ist die systematische Behandlung von Fehlern und die Fähigkeit zur Wiederaufnahme.[Kleppmann 2017, Kap.8]
Teilausfälle können dazu führen, dass ein Job nicht vollständig verarbeitet wird oder dass unklar ist, ob ein Verarbeitungsschritt
bereits erfolgreich abgeschlossen wurde.[Kleppmann 2017, Kap.8]
In der Praxis wird deshalb häufig mit Wiederholungsmechanismen gearbeitet, insbesondere um transiente Fehler (z.B. kurzzeitige Nichterreichbarkeit
einer Datenbank oder eines externen Dienstes) zu überbrücken.[Nygard 2018, Kap.5]
Dabei sind Retry-Strategien mit Backoff wichtig, um durch zeitlich gestaffelte Wiederholungen eine Überlastung des Systems zu vermeiden und die Stabilität zu erhöhen.[Nygard 2018, Kap.5]
Ergänzend ist eine Fehlerklassifikation relevant, da sie bestimmt, ob ein erneuter Versuch sinnvoll ist oder ob ein Job kontrolliert in einen
definierten Fehlerzustand überführt werden sollte.[Nygard 2018, Kap.5]
Wiederaufnahme bedeutet in diesem Kontext, dass Jobs nach einem Ausfall erneut verarbeitbar werden, ohne dass sie dauerhaft „stecken bleiben“.[Kleppmann 2017, Kap.8–9]
Dazu müssen Systeme Zustandsinformationen so führen, dass sie nach einem Crash konsistent weiterarbeiten können, etwa indem Zwischenzustände erkannt
und Jobs erneut eingeplant werden.[Kleppmann 2017, Kap.8–9]
Gleichzeitig darf die Wiederaufnahme nicht zu inkonsistenten Ergebnissen führen, wenn eine Verarbeitung bereits teilweise erfolgt ist.[Kleppmann 2017, Kap.8]
Damit sind Job-Processing-Systeme typischerweise darauf angewiesen, dass Statusänderungen und fachliche Änderungen über klar definierte Transaktionsgrenzen und konsistente
Zustandsübergänge abgesichert werden.[Kleppmann 2017, Kap.8]


Da zwischen dem Absenden und der tatsächlichen Verarbeitung Zeit vergeht und Teilausfälle auftreten können, müssen Systeme definieren,
wie sie mit Duplikaten, Verzögerungen und Wiederholungen umgehen.[Kleppmann 2017, Kap.8,11] Diese Aspekte sind eng mit der Frage verbunden, wie zuverlässig
ein Auftrag verarbeitet wird und welche Konsistenzzusicherungen dabei gelten.
Daher sind Mechanismen erforderlich um Eindeutigkeit, Status, Wiederaufnahme und Fehlerbehandlung kontrolliert umzusetzen.[Kleppmann 2017, Kap.8–9]

"In principle, a file or database is sufficient to connect producers and consumers: a producer writes
every event that it generates to the datastore, and each consumer periodically polls the datastore
to check for events that have appeared since it last ran. This is essentially what a batch process
does when it processes a day’s worth of data at the end of every day." \parencite[Kap.11 --> Transmitting Event Streams]{kleppmann_designing_2017}

"What happens if nodes crash or temporarily go offline—are any messages lost? As with databases,
durability may require some combination of writing to disk and/or replication (see the sidebar
\href{https://learning.oreilly.com/library/view/designing-data-intensive-applications/9781491903063/ch07.html#sidebar_transactions_durability}{“Replication and Durability”}), which has a cost. If you can afford to sometimes lose
messages, you can probably get higher throughput and lower latency on the same hardware."\parencite[Kap.11 Messaging Systems]{kleppmann_designing_2017}


\section{Konsistenz und konkurrierende Verarbeitung}\label{sec:konsistenz}
\textcite[Kap. 7 --> Transactions]{kleppmann_designing_2017}

\subsection{ACID}\label{subsec:acid}

\subsection{Locking}\label{subsec:locking}
Optimistisches Locking \parencite[Kap.6.2]{richardson_microservices_2018}

\subsection{Transaktionsgrenzen}\label{subsec:transaktionsgrenzen}


\section{Fehlertoleranz und Wiederherstellung}\label{sec:fehlertoleranz}

\href{https://learning.oreilly.com/library/view/designing-data-intensive-applications/9781491903063/ch01.html#sec_introduction_reliability}{“Reliability”}

\subsection{Retry-/Backoff-Strategien}\label{subsec:retry/backoff}

\subsection{Fehlerklassifikationen}\label{subsec:fehlerklassifikationen}

\subsection{Recovery}\label{subsec:recovery}

Dein „remote datastore“ ist praktisch deine PostgreSQL-Jobtabelle (persistenter Job-Store). Das entspricht dem Ansatz „State remote halten“, was Recovery vereinfacht, aber Laufzeit/Contention beeinflusst.
 Dein Worker hat vermutlich lokalen, transienten Zustand während der Verarbeitung (in-memory). Der muss nach Crash nicht „rekonstruiert“ werden, sondern du brauchst Mechanismen wie Lease/Timeout + Retry, um Jobs wieder aufzunehmen.
„Replay state from input“ entspricht bei dir am ehesten: Job erneut ausführen (at-least-once) – und genau deshalb brauchst du Idempotenz, weil „replay“ bei Side-Effects sonst doppelte Ergebnisse erzeugt.
\parencite[Kap.11 Rebuilding a state failure]{kleppmann_designing_2017}

\subsection{Idempotenz}\label{subsec:idempotenz}
Wichtig für exactly-once Semantik \parencite[Kap.~8, Abschnitt\glqq Types Of Delivery Guarantees\grqq]{newman_building_2026}
\parencite[Kap.11 --> Atomic commit revisited \& Idempotence]{kleppmann_designing_2017}
\parencite[Kap.3.3.6]{richardson_microservices_2018}
Message Broker, um doppelte Nachrichten zu Eleminieren \parencite[Kap.3.3.6]{richardson_microservices_2018}
Einen Message Consumer Idempotent machen \parencite[Kap.6.1.6]{richardson_microservices_2018}

\section{Zustandsmodellierung mittels State Machines}\label{sec:zustandsmodellierung}


\section{Stand der Technik}\label{sec:forschungsstand}

Für die praktische Umsetzung existieren unterschiedliche Frameworks, die einzelne Aspekte von Background Job-Processing adressieren.
Quartz Scheduler bietet trigger- bzw. zeitgesteuertes Scheduling und unterstützt Persistenz sowie Clustering zur koordinierten Ausführung.[Terracotta Inc. 2026]
Spring Batch stellt eine Infrastruktur für Batch-Verarbeitung bereit, insbesondere mit Job-/Step-Modell und Persistenz der Ausführungshistorie.[Broadcom Inc. 2026]
JobRunr bietet datenbankgestütztes Background Job-Processing in Java mit integriertem Retry-Handling.[Rosoco BV 2026]
Temporal stellt eine Workflow-Engine für langlebige, robuste Ausführung bereit, erfordert jedoch einen eigenständigen Serverbetrieb
und damit zusätzliche Betriebskomplexität.[Temporal Technologies Inc. 2026]
Diese Lösungen verdeutlichen, dass „Background Job-Processing“ in der Praxis nicht allein durch eine Bibliothek gelöst wird, sondern durch
Architekturentscheidungen zur Persistenz, zur Nebenläufigkeitskontrolle sowie zur Fehlerbehandlung und Wiederaufnahme.[Kleppmann 2017]
Die konkrete Wahl und Kombination dieser Mechanismen hängt von den gewünschten Zusicherungen und den betrachteten Fehlerszenarien ab.[Kleppmann 2017]


